% `tex.web` §1183: two fraction operators in one math list, and TeX will not
% guess which nests inside which.
%
%     if incompleat_noad<>null then
%       begin if c>=delimited_code then
%         begin scan_delimiter(garbage,false); scan_delimiter(garbage,false);
%         end;
%       if c mod delimited_code=above_code then scan_dimen(false,false,false);
%       print_err("Ambiguous; you need another { and }");
%       help3("I'm ignoring this fraction specification, since I don't")
%       ("know whether a construction like `x \over y \over z'")
%       ("means `{x \over y} \over z' or `x \over {y \over z}'.");
%       error;
%       end
%
% The delimiters and the thickness are scanned either way — they have to be, or
% the rest of the formula would read them as material — and then the whole
% specification is *ignored*, list unchanged.
%
% We accepted it and built the fraction. That is not merely a missing message:
% the math list came out a level deeper than TeX's for the rest of the formula,
% which is what made trip.tex's `\mkern` two commands later report `math mode`
% where TeX reports `display math mode`.
%
% Braces are how to say what was meant, and the last case shows it — with them
% there is one fraction per list and nothing to complain about.
\catcode`\{=1 \catcode`\}=2 \catcode`\$=3
\errorcontextlines=5 \hsize=100pt \parindent=0pt
\font\rm=cmr10 \textfont0=\rm \scriptfont0=\rm \scriptscriptfont0=\rm \rm
\font\mi=cmmi10 \textfont1=\mi \scriptfont1=\mi \scriptscriptfont1=\mi
\font\sy=cmsy10 \textfont2=\sy \scriptfont2=\sy \scriptscriptfont2=\sy
\font\ex=cmex10 \textfont3=\ex \scriptfont3=\ex \scriptscriptfont3=\ex
\delcode`(="028300 \delcode`)="029301
\showthe\escapechar
\setbox0\hbox{$A\over B\over C$}
\message{PROBE:two=[\the\wd0]}
\setbox1\hbox{$A\over B\abovewithdelims(( 2pt C$}
\message{PROBE:delims=[\the\wd1]}
\setbox2\hbox{$A\over B\above 3pt C$}
\message{PROBE:above=[\the\wd2]}
\setbox3\hbox{$A\over B$}
\message{PROBE:ok=[\the\wd3]}
\setbox4\hbox{${A\over B}\over C$}
\message{PROBE:braced=[\the\wd4]}
\showthe\escapechar
\end
